\section{Twisted lattices and the product formula}\label{sec:lattices}

\subsection{Integrality at finite places and size at infinity}

Let \(\iota\) denote a prime \(\ell\) together with a field isomorphism
\(\mathbb C\simeq\overline{\mathbb Q_\ell}\). Write
\(|z|_\iota=|\iota(z)|_\ell\), with the normalization \(|\ell|_\ell=\ell^{-1}\),
and let \(\bar\iota\) be the isomorphism obtained by precomposing with complex conjugation.
A twist \(\xi\) is a character of \(U(1)\) of weight \(0\) at the infinite places.
Its set of ramified places \(Q_\xi\) is disjoint from the places above \(MN\)
and from its conjugate. For each \(1\le j\le m\), the power \(\xi^j\) is ramified
at every place of \(Q_\xi\). We say that \(\xi\) is admissible at \(\ell\) if the
norm \(N_\xi\) of its modulus is coprime to \(\ell\).
\leansrc{New/Reduction/Data/LevelAndTwists}{Twists and admissibility}.

\begin{editorialnote}
The index \(\iota\) specifies both the prime and the field isomorphism.
The distinction between \(\iota\) and \(\bar\iota\) above the same prime \(\ell\)
accounts for the complex conjugation in one argument of the Hermitian inner product.
\end{editorialnote}

Let \(W_\xi\) be the space of forms into which \(\Theta_\xi\) maps, and let
\(W^{\mathrm{int}}_{\xi,\iota}\) be its module of forms integral at \(\iota\).
Integrality means that the \(\ell\)-adic Weyl gauge of every value of the form
is at most \(1\). For \(\xi\) admissible at \(\ell\), we have
\[
  \phi\in W^{\mathrm{int}}_{\xi,\bar\iota},\quad
  \psi\in W^{\mathrm{int}}_{\xi,\iota}
  \quad\Longrightarrow\quad
  |\langle\phi,\psi\rangle|_\iota\le1 .
\]
Indeed, the inner product is the sum of the Weyl pairings at class
representatives. At this level the stabilizers of the representatives are
trivial, so their orders contribute no denominators. At each point of
evaluation, write \(A\) for the local matrix, \(\rho_{\mathrm W}\) for the
Weyl representation, and \(G\) for the integral Gram matrix. After conjugating
the first argument, apply \(\rho_{\mathrm W}(A^{\mathsf T})\) to it and
\(\rho_{\mathrm W}(A^{-1})\) to the second argument. The gauge bounds imply that
both resulting vectors have integral coordinates. Since
\(\rho_{\mathrm W}(A^{\mathsf T})^{\mathsf T}G\rho_{\mathrm W}(A^{-1})=G\),
the pairing computed in these integral coordinates equals the
\(\iota\)-image of the original pairing. \leanref{gramForm_mem_embValRing}
For the trivial twist, take as coordinates the values at class
representatives in the standard Weyl basis. The inner product in these
coordinates has a positive-definite matrix with rational entries.
\leanref{Data.targetLattice}

\begin{definition}
A vector \(y\in V\) is integral at \(\iota\) if \(\Theta_\xi y\) lies in
\(W^{\mathrm{int}}_{\xi,\iota}\) for every \(\xi\) admissible at \(\ell\).
A finite family \(y=(y_1,\ldots,y_n)\) has logarithmic index at most \(\delta\)
if there are nonnegative reals \(\delta_\ell\) with the following properties.
\begin{enumerate}
\item For every finite set \(P\) of primes,
\(\sum_{\ell\in P}\delta_\ell\log\ell\le\delta\).
\item For each \(\iota\) above \(\ell\), there is an \(n\times n\) matrix
\(U_\iota\) with entries integral at \(\iota\) such that every component of
\(yU_\iota\) is integral at \(\iota\) and
\(|\det U_\iota|_\iota\ge\ell^{-\delta_\ell}\).
\end{enumerate}
The same \(U_\iota\) serves for every twist admissible at that prime.
\end{definition}

Apply \(\sigma\in\operatorname{Aut}(\mathbb C)\) to the coordinates of the
untwisted lift, and write \(\|\Theta_1 y_i\|_\sigma^2\) for the squared norm
computed using the rational inner product matrix. If the Gram determinant
of the lifted family \(y\), computed with these coordinates and this inner
product, is at most \(e^c\) for every \(\sigma\), we say that the logarithm of
the squared covolume is at most \(c\). In particular, if
\(\|\Theta_1 y_i\|_\sigma^2\le e^B\) for every \(i,\sigma\), then Hadamard's
inequality gives the bound \(c=nB\). \leanref{IndexLe}

\subsection{The product formula for heights}

\begin{definition}
A family \(u=(u_\iota)_\iota\) has height at most \(C\), written
\(\operatorname{ht}(u)\le C\), if there are a common multiplier
\(H\in\mathbb C^\times\), a real number \(\alpha\), and nonnegative reals
\(e_\ell\) such that for every \(\sigma,\iota\) and every finite set \(P\)
of primes,
\[
 |\sigma H|\le e^\alpha,\qquad
 |H u_\iota|_\iota\le\ell^{e_\ell},\qquad
 \alpha+\sum_{\ell\in P} e_\ell\log\ell\le C .
\]
\end{definition}

When multiplying two families, we can multiply their common multipliers
and add their \(\alpha\)'s and \(e_\ell\)'s. Thus height bounds add under
multiplication. Every upper bound for the height of the constant family
\(1\) is nonnegative. To see this, take a multiplier \(H\) for that family.
Its absolute values \(|H|_\iota\) are bounded uniformly in \(\iota\) above
\(\ell=2\), so \(H\) is algebraic. In fact, a transcendental \(H\) can be
mapped by a suitable isomorphism to any transcendental element of
\(\overline{\mathbb Q_2}\). The elements algebraic over \(\mathbb Q\)
are countable, and their complement is dense, so there are transcendental
elements of arbitrarily large absolute value. The extension of the
isomorphism is the one used in the algebraicity criterion.
\leanref{isAlgebraic_of_norm_complexToCp_le}

Set \(Q=\operatorname{Nm}_{\mathbb Q(H)/\mathbb Q}H\) and
\(n_H=[\mathbb Q(H):\mathbb Q]\). Then \(Q\ne0\). Writing \(Q\) as a product
of conjugates of \(H\) gives \(\log|Q|\le n_H\alpha\) at the infinite place
and \(\log|Q|_\ell\le n_H e_\ell\log\ell\) at the finite places. Take \(P\)
to contain the primes dividing the numerator or denominator of \(Q\).
The product formula over \(\mathbb Q\) gives
\[
 0=\log|Q|+\sum_{\ell\in P}\log|Q|_\ell
   \le n_H\left(\alpha+\sum_{\ell\in P} e_\ell\log\ell\right)
   \le n_H C .
\]
Thus \(C\ge0\). In particular, if the constant family \(z\ne0\) has height
at most \(C\) and \(|\sigma z|\le e^{C'}\) for every \(\sigma\), then
\(C+C'\ge0\). Indeed, the constant family \(z^{-1}\) has height at most
\(C'\) with multiplier \(H=z\), and the two families have product \(1\).
\leanref{HeightLe.nonneg_of_one}

\subsection{The Gram determinant of a basis}

Assume that the maps \(\Theta_\xi\) are injective and put
\(Z(\xi)=\det(\Theta_\xi^*\Theta_\xi)\). For each \(\iota\), choose a
twist \(\xi_\iota\) admissible at its prime. Let \(y=(y_1,\ldots,y_n)\)
be linearly independent, and complete it by \(z=(z_1,\ldots,z_{n'})\)
to a basis of \(V\). Write \(w=(y,z)\) and set
\[
 G_\xi(w)=\bigl(\langle\Theta_\xi w_i,\Theta_\xi w_j\rangle\bigr)_{ij},
 \qquad \Delta_\xi(w)=\det G_\xi(w).
\]

\begin{proposition}\label{prop:gram-product}
Suppose that the logarithm of the squared covolume of \(y\) is at most
\(c\) and that its logarithmic index is at most \(\delta\). For each twist,
suppose there is a scalar \(\gamma(\xi)\ne0\) such that
\(\gamma(\xi)\Theta_\xi z_j\) is integral at every place where the twist is
admissible. Suppose also that \(\|\Theta_1 z_j\|_\sigma^2\le e^{B_{\mathrm{Ar}}}\)
for all \(\sigma,j\). If the families \((|\gamma(\xi_\iota)|^{-2})_\iota\)
and \((Z(1)/Z(\xi_\iota))_\iota\) have heights at most
\(B_{\mathrm{fin}}\) and \(B_Z\), respectively, then
\begin{equation}\label{eq:gram-product}
  c+2\delta+n'(B_{\mathrm{Ar}}+B_{\mathrm{fin}})+B_Z\ge0.
\end{equation}
\end{proposition}

\begin{proof}
Since \(\bar\iota\) lies above the same \(\ell\) as \(\iota\), both index
matrices \(U_\iota,U_{\bar\iota}\) for \(y\) can be used with
\(\xi=\xi_\iota\). Put
\[
  P_\iota=\operatorname{diag}(U_\iota,\gamma(\xi)I_{n'}),\qquad
  P_{\bar\iota}=\operatorname{diag}(U_{\bar\iota},\gamma(\xi)I_{n'}).
\]
Every entry of \(\overline{P_{\bar\iota}}^{\,t}G_\xi(w)P_\iota\) pairs a form
integral at \(\bar\iota\) with one integral at \(\iota\), so the entry is
integral at \(\iota\). Taking determinants gives
\begin{align*}
 1&\ge|\det P_{\bar\iota}|_{\bar\iota}
             |\Delta_\xi(w)|_\iota\,|\det P_\iota|_\iota,\\
 \bigl|\,|\gamma(\xi)|^{2n'}\Delta_\xi(w)\,\bigr|_\iota
     &\le\ell^{2\delta_\ell}.
\end{align*}
The second inequality uses
\(\overline{\gamma(\xi)}\,\gamma(\xi)=|\gamma(\xi)|^2\) and the lower bounds
on the determinants of the two index matrices. With \(H=1\), the last
inequality gives height at most \(2\delta\) for the family
\((|\gamma(\xi_\iota)|^{2n'}\Delta_{\xi_\iota}(w))_\iota\).

Since \(w\) is a basis of \(V\), the Gram determinant obeys
\(\Delta_\xi(w)=Z(\xi)\det(\langle w_i,w_j\rangle)_{ij}\).
Therefore, for each \(\iota\),
\[
 \Delta_1(w)
 =\bigl(|\gamma(\xi_\iota)|^{2n'}\Delta_{\xi_\iota}(w)\bigr)
     \bigl(|\gamma(\xi_\iota)|^{-2}\bigr)^{n'}
     \frac{Z(1)}{Z(\xi_\iota)} .
\]
It follows that the constant family \(\Delta_1(w)\) has height at most
\(2\delta+n'B_{\mathrm{fin}}+B_Z\).

For the Archimedean bound, apply Fischer's inequality to the Gram
determinant of the full basis after applying \(\sigma\) to the coordinates.
It is at most the product of the determinants of the diagonal blocks for
\(y\) and \(z\). Hadamard's inequality on the auxiliary block for \(z\)
then bounds the determinant of the whole basis by \(e^{c+n'B_{\mathrm{Ar}}}\).
To bound the conjugate \(\sigma(\Delta_1(w))\) itself for an arbitrary
field automorphism \(\sigma\), we use exterior products to account for
complex conjugation. Put \(\sigma^\dagger(z)=\overline{\sigma(\bar z)}\).
The inner product matrix is rational, so \(\sigma(\Delta_1(w))\) is the
inner product of the two exterior products formed using \(\sigma^\dagger\)
and \(\sigma\), respectively. Each has squared norm at most
\(e^{c+n'B_{\mathrm{Ar}}}\). The Cauchy--Schwarz inequality gives
\(\left|\sigma(\Delta_1(w))\right|\le e^{c+n'B_{\mathrm{Ar}}}\).

Injectivity of \(\Theta_1\) gives \(\Delta_1(w)\ne0\). Combining its height
and Archimedean bounds by the product formula proves \eqref{eq:gram-product}.
\leanref{covolume_add_index_nonneg}
\end{proof}
